\providecommand{\todo}[1]{\textbf{[TODO: #1]}}

บทนี้ทบทวนงานวิจัยและเทคโนโลยีที่เกี่ยวข้องกับการแปลงเสียงเป็นข้อความในงานทางการแพทย์ เพื่อชี้ให้เห็นช่องว่างที่โครงงานนี้ต่อยอด

\section{งานวิจัยด้านการรู้จำเสียงในงานทางการแพทย์}

\subsection{Whisper}
\begin{itemize}
    \item ผู้วิจัย: Radford และคณะ (เผยแพร่ครั้งแรกปี 2022 ตีพิมพ์ใน ICML 2023) \cite{radford2023whisper}
    \item แนวคิดหลัก: ฝึกแบบจำลองให้ทำนายข้อความถอดเสียงจากข้อมูลเสียงบนอินเทอร์เน็ตจำนวนมาก (Weak Supervision) ขยายถึง 680,000 ชั่วโมงแบบหลายภาษาและหลายงาน
    \item วิธีการ: ใช้สถาปัตยกรรม Transformer แบบ Encoder--Decoder สำเร็จรูป มีแบบจำลองหลายขนาด และประเมินแบบ Zero-shot โดยไม่ปรับแต่งกับชุดข้อมูลเฉพาะ
    \item ข้อดีและข้อจำกัด: ผู้เขียนรายงานว่าแบบจำลองใกล้เคียงมนุษย์ในด้านความแม่นยำและความทนทาน แต่งานนี้ไม่ได้ทดสอบกับศัพท์พยาธิวิทยาหรือเสียงตู้ดูดควัน
\end{itemize}

\subsection{ข้อผิดพลาดของการรู้จำเสียงในเอกสารทางคลินิก}
\begin{itemize}
    \item ผู้วิจัย: Zhou และคณะ (2018) \cite{zhou2018errors}
    \item แนวคิดหลัก: วัดข้อผิดพลาดของเอกสารทางคลินิกที่เกิดจากการบอกให้ซอฟต์แวร์รู้จำเสียงช่วยเขียนในสามขั้นตอน ได้แก่ ผลดิบจากซอฟต์แวร์ ฉบับที่ผู้ถอดความมืออาชีพแก้ไข และฉบับสุดท้ายที่แพทย์ลงนาม
    \item วิธีการ: สุ่มเอกสาร 217 ฉบับจากระบบสุขภาพ 2 แห่ง แล้วจัดประเภทและนับข้อผิดพลาดทุกขั้นตอน
    \item ผลลัพธ์: อัตราข้อผิดพลาดเท่ากับ 7.4\% ในผลดิบ 0.4\% หลังผู้ถอดความตรวจแก้ และ 0.3\% ในฉบับที่แพทย์ลงนาม และข้อผิดพลาดจำนวนมากเกี่ยวข้องกับข้อมูลทางคลินิก
    \item ข้อดีและข้อจำกัด: สนับสนุนว่าต้องมีขั้นตอนตรวจทานโดยมนุษย์ซึ่งโครงงานนี้ออกแบบไว้ แต่เป็นการบอกให้เขียนแบบเบื้องหลัง (Back-end Dictation) ด้วยซอฟต์แวร์เชิงพาณิชย์ ไม่ใช่ Whisper และไม่ได้ศึกษาห้องตรวจชิ้นเนื้อ
\end{itemize}

\subsection{การสร้างข้อความที่ไม่มีอยู่ในเสียง (Hallucination)}
\begin{itemize}
    \item ผู้วิจัย: Koenecke และคณะ (2024) \cite{koenecke2024careless}
    \item แนวคิดหลัก: ตรวจว่า Whisper บางครั้งสร้างวลีหรือประโยคที่ไม่มีอยู่ในเสียงต้นฉบับ
    \item วิธีการ: ส่งเสียง 13,140 ช่วงจาก AphasiaBank (ผู้พูด 437 คน ทั้งผู้มีภาวะ Aphasia และกลุ่มควบคุม) เข้า Whisper API ในปี 2023 และจัดประเภทข้อความที่สร้างเกินเข้ามา
    \item ผลลัพธ์: ประมาณ 1.4\% ของช่วงเสียงมีข้อความที่สร้างเกิน ร้อยละ 38 ของกลุ่มนั้นมีเนื้อหาที่เป็นอันตราย เช่น ชื่อหรือสถานะสุขภาพที่แต่งขึ้น และพบมากขึ้นในเสียงที่มีสัดส่วนช่วงไม่มีเสียงพูดยาว ผู้เขียนรายงานว่าระบบของผู้ให้บริการรายอื่นที่ทดสอบไม่พบกรณีเทียบเท่าในชุดเสียง 187 ช่วงที่ก่อปัญหา
    \item ข้อดีและข้อจำกัด: ชี้ความเสี่ยงของการใช้ Whisper ในงานทางการแพทย์ และสอดคล้องกับการตัดช่วงเงียบก่อนถอดความและการตรวจทานผลโดยมนุษย์ในโครงงานนี้ (เป็นข้อสังเกตของผู้จัดทำ ไม่ใช่ข้อสรุปของงานต้นฉบับ) งานนี้ทดสอบ Whisper API กับเสียงผู้พูดที่มีภาวะ Aphasia ไม่ใช่รุ่น Small แบบออฟไลน์
\end{itemize}

\section{งานวิจัยด้านเทคนิคหรือเครื่องมือที่เกี่ยวข้อง}

\subsection{เทคนิคที่ใช้}
\begin{itemize}
    \item \textbf{Quantization แบบ INT8:} แปลงน้ำหนักของแบบจำลองเป็นจำนวนเต็ม 8 บิตเพื่อลดหน่วยความจำและเวลาประมวลผล งานต้นแบบของ Jacob และคณะรายงานว่าลดหน่วยความจำได้ราว 4 เท่าเมื่อเทียบกับเลขทศนิยม 32 บิตโดยความแม่นยำใกล้เคียงเดิม แต่ทดสอบกับแบบจำลองด้านภาพ ไม่ใช่แบบจำลองเสียง \cite{jacob2018quantization}
    \item \textbf{Spectral Denoising (afftdn):} ตัวกรองใน FFmpeg ที่ประมาณและลดทอนเสียงรบกวนในโดเมนความถี่ มีค่าเริ่มต้น nr 12 dB, nf $-50$ dB และชนิดเสียงรบกวนแบบขาว \cite{ffmpeg_afftdn}
    \item \textbf{การตัดช่วงเงียบ (Silence Removal):} ตัดช่วงที่ระดับสัญญาณต่ำกว่าเกณฑ์ออกเพื่อลดเวลาประมวลผลและความเสี่ยงของ Hallucination โครงงานนี้ใช้ตัวกรอง silenceremove ของ FFmpeg ซึ่งตัดตามระดับสัญญาณ ไม่ใช่แบบจำลองตรวจจับเสียงพูด (VAD)
    \item \textbf{Prompt Conditioning:} ป้อนข้อความนำเพื่อชี้นำการสะกดคำศัพท์เฉพาะทาง Whisper รองรับการปรับแต่งนี้ \cite{radford2023whisper} และเอกสารของ OpenAI (ที่ Koenecke และคณะอ้างถึง \cite{koenecke2024careless}) ระบุว่าใช้แก้คำหรือคำย่อที่แบบจำลองมักถอดผิดได้
    \item \textbf{การสกัดข้อมูลแบบกำหนดแน่นอน:} ใช้กฎและนิพจน์ปรกติแทนแบบจำลองเชิงกำเนิด ทำให้ผลตรวจสอบย้อนกลับได้และไม่สร้างค่าที่ไม่ปรากฏในข้อความ
    \item \textbf{โปรโตคอล CAP และระบบ pTNM:} CAP กำหนดข้อมูลที่ต้องรายงานสำหรับมะเร็งเต้านมชนิดลุกลาม โดยมีโปรโตคอลแยกสำหรับ DCIS และอ้างอิงเกณฑ์ pTNM ของ AJCC \cite{cap_breast,ajcc2017}
\end{itemize}

\subsection{เครื่องมือและซอฟต์แวร์}
\begin{itemize}
    \item \textbf{CTranslate2:} เอนจินอนุมานสำหรับแบบจำลอง Transformer ที่รองรับ Whisper และการทำ Quantization บน CPU โดยผู้พัฒนารายงานว่าลดขนาดแบบจำลองได้ราว 4 เท่าโดยความแม่นยำลดลงน้อย \cite{ctranslate2}
    \item \textbf{Faster-Whisper:} การนำ Whisper มาเขียนใหม่บน CTranslate2 ผู้พัฒนารายงานว่าเร็วกว่า openai/whisper สูงสุดราว 4 เท่าที่ความแม่นยำเท่ากัน และใช้หน่วยความจำน้อยกว่า ตัวเลขนี้มาจากเอกสารของผู้พัฒนา ไม่ใช่การประเมินอิสระ \cite{fasterwhisper}
    \item \textbf{Vosk (Kaldi):} เครื่องมือรู้จำเสียงแบบออฟไลน์ที่สร้างบน Kaldi ใช้แบบจำลองขนาดเล็กราว 50 MB ใช้เป็นแบบจำลองเปรียบเทียบ \cite{vosk,povey2011kaldi}
    \item \textbf{wav2vec 2.0:} แบบจำลองเรียนรู้เสียงแบบ Self-supervised แล้วปรับแต่งด้วยข้อมูลถอดเสียงจำนวนน้อย ใช้เป็นแบบจำลองเปรียบเทียบ \cite{baevski2020wav2vec}
    \item \textbf{FFmpeg, Python, Flask, PostgreSQL, SQLite:} ใช้ประมวลผลเสียง พัฒนาระบบ และจัดเก็บข้อมูล
\end{itemize}

\section{การวิเคราะห์ช่องว่างของงานวิจัย}

จากการทบทวน พบว่า Whisper ทำงานออฟไลน์ได้แต่ไม่ได้ปรับให้เข้ากับเสียงตู้ดูดควันและศัพท์เฉพาะ งานด้านเอกสารทางคลินิกแสดงว่าการรู้จำเสียงยังต้องมีการตรวจทาน และงานด้านความปลอดภัยชี้ความเสี่ยงของการสร้างข้อความเกินจริง ขณะที่ระบบถอดความทั่วไปมักหยุดที่ข้อความดิบและไม่แปลงเป็นฟิลด์ โครงงานนี้จึงรวมการกรองเสียง การถอดความแบบออฟไลน์ที่เร็วขึ้น และการสกัดฟิลด์แบบกำหนดแน่นอนไว้ในระบบเดียว ดังสรุปในตารางที่ \ref{tab:gap}

\begin{table}[H]
\centering
\caption{การเปรียบเทียบช่องว่างของงานวิจัย}
\label{tab:gap}
\small
\begin{tabular}{|p{4.0cm}|p{5.0cm}|p{5.0cm}|}
\hline
\textbf{มิติ} & \textbf{Whisper Small (baseline)} & \textbf{PathoWhisper} \\
\hline
การทำงานออฟไลน์ & ออฟไลน์ได้ แต่ช้าบน CPU (20.04 วินาที/กรณี) & ออฟไลน์ทั้งหมด (11.50 วินาที/กรณี) \\
\hline
เสียงตู้ดูดควัน ($-10$ dB) & WER 44.88\% & WER 31.94\% \\
\hline
การพูดแก้ไขกลางประโยค & ถอดตามเสียง ไม่มีขั้นสกัดค่า & ตัวสกัดเลือกค่าที่พูดทีหลัง \\
\hline
การกรอกฟิลด์การบรรยาย & ได้ข้อความดิบ & สกัด 15 ฟิลด์และสร้าง PDF \\
\hline
\end{tabular}
\end{table}

\noindent บริการรู้จำเสียงบนคลาวด์ไม่ได้ถูกนำมาทดสอบในโครงงานนี้ จึงไม่รวมในตารางข้างต้น ข้อจำกัดสำคัญของบริการลักษณะนี้ในงานที่มีข้อมูลสุขภาพคือเสียงต้องออกนอกเครื่องไปประมวลผล ซึ่งเป็นเหตุผลด้านความเป็นส่วนตัวที่โครงงานนี้เลือกทำงานแบบออฟไลน์